\newpage
\section{第25章：System F 的 ML 实现}

\begin{taplauthorsolutionentrybox}
\phantomsection\label{sol:author:25.2.1}
\textbf{25.2.1 解答}

记 \(\uparrow_c^d(T)\) 为在截止位置 \(c\) 之上把自由类型变量上移 \(d\)
位：
\[
\begin{array}{rcl}
\uparrow_c^d(k)&=&
 \begin{cases}
 k+d&k\geq c,\\
 k&k<c,
 \end{cases}\\
\uparrow_c^d(T_1\to T_2)&=&
 \uparrow_c^d(T_1)\to\uparrow_c^d(T_2),\\
\uparrow_c^d(\forall X.T_2)&=&
 \forall X.\uparrow_{c+1}^d(T_2),\\
\uparrow_c^d(\exists X.T_2)&=&
 \exists X.\uparrow_{c+1}^d(T_2).
\end{array}
\]
顶层移位简写为 \(\uparrow^d(T)=\uparrow_0^d(T)\)。
\taplsolutionbacklink{ex:25.2.1}{返回练习25.2.1}
\end{taplauthorsolutionentrybox}

\translatornote{在书中的无名运行时表示里，每个变量节点还保存了该位置预期的
上下文长度；移位时，这个长度也要同步增加 \(d\)。}

\begin{taplauthorsolutionentrybox}
\phantomsection\label{sol:author:25.4.1}
\textbf{25.4.1 解答}

这次移位是为类型变量 \(X\) 腾出上下文位置。把 \(v_{12}\) 代入 \(t_2\)
所得的结果应当在 \(\Gamma,X\) 中具有良好的作用域，而原来的 \(v_{12}\)
只相对于 \(\Gamma\) 定义。
\taplsolutionbacklink{ex:25.4.1}{返回练习25.4.1}
\end{taplauthorsolutionentrybox}

\translatornote{具体顺序是：先把 \(v_{12}\) 上移一位，再替换项变量；随后
\texttt{tytermSubstTop} 代入实际类型并删除 \(X\) 的上下文位置，使最终结果
重新相对于 \(\Gamma\) 解释。}
